\section{Introduction}

The automated discovery of novel mathematical concepts and theorems has emerged as a fascinating frontier where artificial intelligence meets pure mathematics. Recent advances in machine learning and automated reasoning have opened new possibilities for systematically exploring mathematical structures and identifying promising research directions \cite{halachmi2024ramanujan, tsoukalas2025learning}. However, the fundamental challenge of automatically generating meaningful and mathematically interesting research topics remains largely unsolved.

The process of mathematical discovery has traditionally relied on human intuition, creativity, and deep understanding of existing mathematical structures. While computational tools have long assisted in verification and exploration \cite{subercaseaux2023automated}, the autonomous generation of novel, well-motivated mathematical research directions presents unique challenges. These include formalizing mathematical interestingness, ensuring logical consistency, and maintaining connections to established mathematical theories.

Recent work has demonstrated the potential of automated systems in mathematical discovery, from identifying patterns in integer relations \cite{halachmi2024ramanujan} to generating conjectures in graph theory \cite{davila2025reverie}. However, these approaches typically operate within pre-defined mathematical domains and rely heavily on human-crafted heuristics for determining mathematical significance.

In this paper, we introduce a formal framework for automated research topic generation in mathematics. Our main contributions are:

1. A rigorous formalization of mathematical research topic spaces, including metrics for evaluating novelty, significance, and tractability

2. An algorithmic approach for exploring these spaces using techniques from recursive function theory and functional analysis

3. Theoretical guarantees on the mathematical consistency and relevance of generated research directions

4. A characterization of the relationship between automatically generated topics and classical mathematical structures

Our work builds upon recent advances in automated mathematical reasoning \cite{georgiev2025mathematical} while introducing novel theoretical tools for systematic exploration of the mathematical landscape. The framework we develop provides a foundation for future research in automated mathematical discovery and offers insights into the nature of mathematical creativity itself.

The rest of this paper is organized as follows. Section 2 presents the formal definitions and theoretical foundations. Section 3 develops our main results on topic space exploration. Section 4 provides examples and applications, while Section 5 discusses implications and future directions.